\section{Conclusion}
\label{sec:conclusion}

The results of the study show that AI programming aids are perfectly capable of performing simple to complex tasks and can also provide strong support for developers and newcomers in the field of IaC. Even though there are significant differences in the quality of the output and generated responses, both tools examined demonstrate the ability to solve complex tasks and generate clean, structured and well-founded Terraform, thereby meeting the user's requirements. Nevertheless, Windsurf proves to be a more robust tool that is better suited to this use case in terms of output quality. As a result, it is advisable to use Windsurf when programming IaC with Terraform. However, this study only evaluated Windsurf Cascade and Github Copilot via VSCode and did not consider other systems such as Cursor. The next logical step would therefore be to conduct an expanded study with all major players in the AI programming tools field, using different modes and not limiting ourselves to just ‘edit’ or ‘chat’ mode, but also modes such as agent mode, which is becoming increasingly popular. The tools should also be tested with difficult tasks using existing codebases in order to examine how successfully the tools adapt to existing coding styles.